\scanpage{262}
- \mspage{94} - \margin{IX}

la biextension \(W^{\mathrm{ab}}\) de \((B,B')\) par \(\mathbb G_{mS}\) est une
dualité parfaite sur \(U\) (ou ce qui revient au même si \(U\) est dense dans
\(S\), sur \(\unsure{S}\)). C'est une propriété qu'il suffit de vérifier en tout
point \(\eta\) de \(\unsure{U}\), et alors l'argument habituel nous ramène au cas où
\(S\) est un \unsure{trait} de point générique \(\eta\), qui est déjà connu.

La relation d'orthogonalité nous permet d'autre part d'étendre au cas présent
les formules (5.5.7), (5.5.8) et (5.5.9), où on remplace l'opération de restriction
au point générique \(\eta\) par celle de restriction à l'ouvert \(U\) ; par suite,
5.6 s'étend également au cas présent. Nous nous dispensons de réécrire ici les
formules et l'énoncé précédents dans le contexte présent.

Résumons, du point de vue des modèles de Néron, les résultats obtenus concernant
les parties abéliennes et leurs accouplements :

\uline{Proposition} 7.5. \uline{Soit} \(S\) \uline{un trait complet},
\uline{de corps de fonctions} \(K\). A \uline{tout schéma abélien} \(A_K\)
\uline{sur} \(K\) \uline{ayant une réduction semi-stable sur} \(S\),
\uline{le procédé de} 7.1,\uline{appliqué à la composante neutre du modèle de Néron}
\(A\), \uline{de} \(A_K\) \uline{permet de définir la} partie abélienne
\(B=A^{\mathrm{ab}}\) \uline{de celui-ci}, \uline{qui}
\uline{est un schéma abélien sur} \(S\), \uline{dépendant fonctoriellement de}
\(A_K\). \uline{Pour} \uline{tout autre schéma abélien} \(A'_K\)
\uline{à réduction semi-stable}, \uline{on peut à toute}
\uline{correspondance divisorielle} \(W_K\) \uline{sur} \((A_K,A'_K)\),
\uline{associer une correspondance} \uline{divisorielle unique} \(W^{\mathrm{ab}}\)
\uline{sur} \((A^{\mathrm{ab}},A'^{\mathrm{ab}})\),
\uline{telle que pour tout entier} \(n\geq0\), \(W_n^{\mathrm{ab}}\)
\uline{ait comme image inverse par}
\(A_n^0\longrightarrow B_n=A_n^{0\mathrm{ab}}\),
\(A'^0_n\longrightarrow B'_n=A'_n{}^{0\mathrm{ab}}\)
\uline{la biextension} \(W_n\) \uline{de} \((A_n^0,A'^0_n)\)
\uline{déduite de} \(W\) \uline{par} \uline{réduction}, \uline{où} \(W\)
\uline{est le prolongement canonique} (1.4) \uline{de} \(W_K\) \uline{à}
\((A^0,A'^0)\).
\margin{\figure{\begin{tikzpicture}[scale=.3]
\draw (0,.4)--(.2,.4);\fill (.15,0) circle (.035);\fill (.3,-.06) circle (.035);
\end{tikzpicture}}}
\margin{412}
\margin{412}
\margin{261}
